\section{Introducción}
La oximetría de pulso es una técnica no invasiva de amplia utilización en la instrumentación biomédica para la monitorización continua y en tiempo real de la saturación periférica de oxígeno en sangre ($SpO_2$) y la frecuencia cardíaca. Esta tecnología se basa en la combinación de la espectrofotometría de absorción óptica y la fotopletismografía para registrar las variaciones volumétricas del flujo sanguíneo en la vascularización periférica durante el ciclo cardíaco.

Los objetivos de esta practica comprenden:
\begin{itemize}
    \item Familiarizarse con los distintos formatos, geometrías y características constructivas de los sensores de oximetría de pulso ($SpO_2$) comerciales.
    \item Identificar la distribución de pines en el conector DB9 del sensor BCI mediante mediciones de resistencia y continuidad, reconociendo las terminales de los emisores (LEDs rojo e infrarrojo), del fotodetector y del blindaje del cable paciente.
    \item Analizar e identificar las distintas etapas de acondicionamiento analógico y digital presentes en los circuitos comerciales e instrumentales de saturación de oxígeno.
    \item Implementar en protoboard un circuito de sensado de intensidad de luz para los emisores rojo e infrarrojo utilizando el sensor BCI estándar y evaluar los niveles de respuesta obtenidos.
    \item Evaluar el comportamiento de un circuito de acondicionamiento de señal pletismográfica y discutir el diseño de una segunda etapa de amplificación y filtrado para optimizar la calidad de la señal medida.
\end{itemize}
